\documentclass[12pt,a4paper]{article}
\usepackage[utf8]{inputenc}
\usepackage[spanish,es-noquoting]{babel}
\usepackage[T1]{fontenc}
\usepackage{lmodern}
\usepackage{geometry}
\geometry{top=2cm, bottom=2cm, left=2.5cm, right=2.5cm}
\usepackage{xcolor}
\usepackage{listings}
\usepackage{enumitem}
\usepackage{fancyhdr}
\usepackage{hyperref}
\usepackage{tikz}
\usetikzlibrary{arrows.meta, positioning, calc, shapes.geometric}

\definecolor{codegreen}{rgb}{0,0.6,0}
\definecolor{backcolour}{rgb}{0.95,0.95,0.92}

\lstset{
    language=C,
    backgroundcolor=\color{backcolour},
    commentstyle=\color{codegreen},
    keywordstyle=\color{blue},
    basicstyle=\ttfamily\small,
    breaklines=true,
    frame=single,
    numbers=left,
    tabsize=4,
    extendedchars=true,
    showstringspaces=false,
    literate={á}{{\'a}}1 {é}{{\'e}}1 {í}{{\'i}}1 {ó}{{\'o}}1 {ú}{{\'u}}1 {ñ}{{\~n}}1
}

\pagestyle{fancy}
\fancyhf{}
\lhead{Monitorías Sistemas Operativos 2026-I}
\rhead{Pre-Taller: Señales}
\cfoot{\thepage}

\title{\textbf{Pre-Taller 1: Protocolo de Telemetría por Señales}}
\author{Malak Sanchez \\ \small Monitorías de Sistemas Operativos}
\date{\today}

\begin{document}

\maketitle

\section*{Problema: Transmisión de Coordenadas de Rastreo (GPS)}
Se requiere implementar un sistema de telemetría donde un proceso \textbf{hijo} simula ser un rastreador GPS y un proceso \textbf{padre} actúa como el servidor de monitoreo.

El hijo debe leer un archivo de texto (\texttt{posiciones.txt}) que contiene pares de números enteros representando coordenadas cartesianas $(X, Y)$ por línea. Su misión es transmitir estas coordenadas al padre utilizando \textbf{únicamente señales}.

\subsection*{Especificaciones del Protocolo}
Dado que las señales no permiten pasar datos directamente, usted debe proponer e implementar un mecanismo de codificación. Algunas estrategias sugeridas son:
\begin{itemize}
    \item \textbf{Modulación por Ancho de Pulsos (Conteo):} Enviar $N$ señales para representar el valor de $X$, hacer el cambio de eje, y $M$ señales para $Y$.
\end{itemize}

\begin{center}
\begin{tikzpicture}[node distance=3cm, font=\small, >=Stealth]
    \node[draw, rounded corners, fill=orange!10, minimum width=2.8cm, minimum height=1.2cm, align=center] (hijo) {HIJO \\ (Rastreador GPS)};
    \node[draw, rounded corners, fill=cyan!10, minimum width=2.8cm, minimum height=1.2cm, align=center, right=5cm of hijo] (padre) {PADRE \\ (Servidor de Monitoreo)};
    
    \draw[->, thick, dashed] (hijo.east) -- node[midway, above] {Ráfagas de Señales} node[midway, below] {SIGUSR1 / SIGUSR2} (padre.west);
\end{tikzpicture}
\end{center}

\textbf{Ejemplo de ejecución esperado:} 
\begin{lstlisting}[numbers=none]
$ ./gps_telemetria posiciones.txt
[Padre] Servidor listo. Esperando datos del rastreador...
[Padre] Coordenada recibida: X=12, Y=8
[Padre] Coordenada recibida: X=45, Y=102
[Padre] Fin de trayectoria detectado.
\end{lstlisting}

\section*{Requerimientos Técnicos}
\begin{enumerate}
    \item El hijo debe manejar los tiempos de envío cuidadosamente para evitar que el padre pierda señales.
    \item Está estrictamente \textbf{prohibido} el uso de tuberías, archivos temporales o memoria compartida para el intercambio de las coordenadas.
\end{enumerate}

\end{document}
